% Part: sets-functions-relations
% Chapter: induction
% Section: induction

\documentclass[../../include/open-logic-section]{subfiles}

\begin{document}

\olfileid[pt]{sfr}{ind}{int}

\olsection{Introdução}

\begin{explain}
A indução é uma técnica de demonstração muito usada na matemática e na lógica.
Talvez você já a conheça da matemática, na qual se faz indução sobre os
números naturais. Na verdade, há muitos tipos diferentes de conjuntos de
objetos matemáticos e lógicos aos quais a indução se aplica.

Em geral, a indução permite demonstrar uma afirmação universal, isto é,
mostrar que todo objeto de certo tipo tem alguma propriedade. Em particular,
a indução costuma ser útil quando um método de demonstração por ``introdução
universal'' é complicado demais. A indução só funciona sobre objetos
matemáticos construídos de maneiras especiais: se cada elemento do conjunto
é básico ou é construído a partir de elementos básicos, podemos usar indução
sobre ele. Mais precisamente:
\end{explain}

\begin{defn}
Um conjunto $S$ é \emph{fechado} sob uma função $f$ se, e somente se,
sempre que $x \in S$, então $f(x) \in S$.
\end{defn}

\begin{explain}
Por exemplo, os números naturais ($\mathbb{N}$) são fechados sob a função
sucessor $f(x) = x+1$. Se $x$ é um número natural, $x+1$ também é. No
entanto, os números naturais não são fechados sob a função predecessor
$f(x) = x-1$, pois $0$ é um número natural, mas $-1$ não é.
\end{explain}

\begin{defn}
Diz-se que uma propriedade $P$ é \emph{preservada por} uma função $f(x)$
quando, para qualquer objeto $a$ no domínio de $f$, $P(a)$ implica
$P(f(a))$ (se $a$ tem a propriedade $P$, então $f(a)$ também a tem).
\end{defn}

\begin{explain}
A propriedade de ``ser um número par'' é preservada pela função
$f(x) = x+2$, pois, se $x$ é par, $x+2$ também é. A propriedade de ser
um número par não é preservada pela função sucessor, pois, se $x$ é par,
$x+1$ é ímpar.

A indução funciona sobre os números naturais porque todo número natural
pode ser alcançado a partir de 0 aplicando-se a função sucessor um número
finito de vezes. Essa é uma característica de qualquer conjunto sobre o
qual podemos fazer indução.
\end{explain}

\begin{defn}[Indução sobre $\Nat$]
Se 0 tem alguma propriedade $P$, e se $P$ é preservada pela função
sucessor, então todo número natural tem a propriedade $P$.
\end{defn}

\begin{ex} A soma dos primeiros $n$ números naturais é $\frac{n(n+1)}{2}$. \\

\textbf{Caso base.} A soma dos primeiros 0 números naturais é $0=
\frac{0\cdot 1}{2}$.\\

\textbf{Passo indutivo.} Suponha que a propriedade valha para $k$.
Mostraremos que ela vale para $k+1$. Em outras palavras, suponha que
$P(k)$ valha, isto é,
\[ 0 + 1 + \ldots + k = \frac{k(k+1)}{2} \]
Mostraremos que $P(k+1)$ vale, O.s.\@, que
\[ 0 + 1 + \ldots + k + (k+1) = \frac{(k+1)(k+2)}{2} \]
usando a hipótese de que $P(k)$.

\begin{align*}
0 + 1 + \ldots + k &= \frac{k(k+1)}{2} \tag{Hipótese} \\
0 + 1 + \ldots + k + (k+1) &= \frac{k(k+1)}{2} + (k+1)
\tag{Somar $k+1$ aos dois lados} \\
&= \frac{k(k+1) + 2(k+1)}{2} \\
&= \frac{2k + k + 2k +2}{2} \\
&=\frac{(k+1)(k+2)}{2}
\end{align*}
Portanto, o passo indutivo é válido, e demonstramos a afirmação.
\end{ex}

\begin{explain}
Para !!{formula}s, os elementos básicos são as !!{formula}s atômicas.
As !!{formula}s mais complexas são obtidas reunindo-se !!{formula}s menos
complexas por meio de operadores. Pode-se considerar que cada operador
corresponde a uma função de !!{formula}s que devolve uma nova !!{formula},
reunindo as duas com o operador em questão. Por exemplo, a função que reúne
duas !!{formula}s $!A$ e $!B$ em uma conjunção pode ser escrita como
$\mathcal E _\land (!A, !B) = !A \land !B$. Podemos chamá-las de
\emph{funções de construção de fórmulas}. O conjunto das !!{formula}s é
fechado sob essas funções: sempre que $!A$ e $!B$ são !!{formula}s,
$\lnot !A$, $!A \land !B$ etc. também são.
\end{explain}

\begin{defn}[Indução sobre !!{formula}s]
Se toda !!{formula} atômica tem a propriedade $P$, e se $P$ é preservada
pelas funções construtoras de !!{formula}, então toda !!{formula} tem
a propriedade $P$.
\end{defn}

\begin{explain}
Essa definição sugere uma ``receita'' para demonstrações indutivas sobre
!!{formula}s. Se nos pedirem que mostremos que toda !!{formula} tem a
propriedade $P$, seguimos os seguintes passos:\\

\textbf{Caso base.} Seja $!A$ uma !!{formula} atômica. [\ldots]
Portanto, $!A$ tem a propriedade $P$.

\textbf{Passo indutivo.} Sejam $!A$ e $!B$ !!{formula}s que têm a
propriedade $P$.

Caso $\lnot$: [\ldots] Portanto, $\lnot !A$ tem a propriedade $P$.

Caso $\land$: [\ldots] Portanto, $!A \land !B$ tem a propriedade $P$.

Caso $\lor$: [\ldots] Portanto, $!A \lor !B$ tem a propriedade $P$.

Caso $\lif$: [\ldots] Portanto, $!A \lif !B$ tem a propriedade $P$.

Caso $\liff$: [\ldots] Portanto, $!A \liff !B$ tem a propriedade $P$.

Caso $\lforall[]$: [\ldots] Portanto, $\lforall[x] !A$ tem a propriedade $P$.

Caso $\lexists[]$: [\ldots] Portanto, $\lexists[x] !A$ tem a propriedade $P$. \\

Portanto, toda !!{formula} tem a propriedade $P$.
\end{explain}

\end{document}
